% =====================================================================
% §1 Moving Target Defence
% =====================================================================

\section{Moving Target Defence}
\label{sec:mtd}

This section introduces the three design questions that organise MTD, traces how the field models, quantifies, and evaluates MTD techniques in practice, and turns to the adaptive-selection frontier where orchestration research now concentrates.



\subsection{MTD Concept and Taxonomy}
\label{subsec:mtd_concept}

MTD is the timely manipulation of system configurations to modify and control the \emph{attack surface}---the set of points at which an attacker engages with the system~\cite{cho2020}.

The MTD design space is organised around three questions: \emph{what} to move, \emph{how} to move it, and \emph{when} to move~\cite{cho2020}. The first concerns the configurable attributes available to the defender---IP addresses, ports, network topologies, operating systems, software stacks, virtual machines.

The second is captured by the canonical Shuffle-Diversity-Redundancy (SDR) taxonomy~\cite{cho2020}, whose three primitives are complementary rather than partitioned (Fig.~\ref{fig:sdr-relationships}). \emph{Shuffling} rearranges or randomises existing components---IP mutation, port hopping, topology reconfiguration---invalidating reconnaissance the attacker has already performed. \emph{Diversity} deploys different implementations of the same function, so that an exploit against one variant is unlikely to apply to others. \emph{Redundancy} replicates components to preserve service while the other two operate. Hybrid combinations of SDR primitives are an active research direction~\cite{cho2020}.

The third question---when to move---distinguishes proactive (time-triggered), reactive (event-triggered), and hybrid scheduling regimes~\cite{cho2020}. This is the tension between cost and security: moving too often imposes overhead on legitimate users; moving too rarely leaves the attacker's reconnaissance valid for longer than it should be.

\begin{figure}[t]
\centering
\resizebox{\columnwidth}{!}{%
\begin{tikzpicture}[
  >=stealth,
  every node/.style={font=\small},
  prop/.style={draw, rectangle, minimum height=6mm, inner sep=3pt},
  op/.style={draw, ellipse, minimum width=22mm, minimum height=9mm},
]
  \node[op, fill=green!25]  (div) at (0,2)     {Diversity};
  \node[op, fill=orange!25] (shf) at (-3,-0.6) {Shuffling};
  \node[op, fill=blue!18]   (red) at (3,-0.6)  {Redundancy};

  \node[prop, fill=green!45]  (res)  at (0,3.6)   {Resilience / Robustness};
  \node[prop, fill=orange!35] (perf) at (-3,-2.1) {Performance / Efficiency};
  \node[prop, fill=blue!30]   (rel)  at (3,-2.1)  {Reliability / Availability};

  % property connectors -- plain lines, NO arrowheads
  \draw (res) -- (div);
  \draw (shf) -- (perf);
  \draw (red) -- (rel);

  % SDR relationships -- bidirectional, matching Cho Fig. 4
  \draw[<->] (shf) to[bend left=25]  (div);
  \draw[<->] (div) to[bend left=25]  (red);
  \draw[<->] (shf) to[bend right=25] (red);
\end{tikzpicture}%
}
\caption{Relationships between Shuffling, Diversity, and Redundancy (SDR). Adapted from \cite[Fig.~4]{cho2020}.}
\label{fig:sdr-relationships}
\end{figure}



\subsection{MTD Evaluation Paradigms}
\label{subsec:mtd_eval}

An MTD evaluation study selects a network setting --- enterprise, cloud,
software-defined, Internet of Things (IoT), or cyber-physical --- adopts a
threat model, deploys one or more MTD mechanisms, and assesses the result
against a metric set using one of several validation methods~\cite{cho2020}.
This section examines the modelling traditions, metric families, and
validation methods these choices draw from.

Modelling techniques cluster into game-theoretic, evolutionary, and
learning-based traditions, with game-theoretic formulations dominating by
volume~\cite{cho2020}. A representative game-theoretic formulation casts the
defender as a leader committing to a configuration-switching strategy and the
attacker as a follower that performs reconnaissance before launching, typically
drawing on exploits from public enumerations such as the Common Vulnerabilities
and Exposures (CVE) database, with attacker-type uncertainty encoded as a prior
distribution~\cite{cho2020}.

Cho et al.~\cite{cho2020} partition the MTD metric space along two axes
(Table~\ref{tab:mtd-metrics}): \emph{perspective} (attacker-side or
defender-side) and \emph{purpose} (effectiveness, capturing how well the MTD
reduces attack success, and efficiency, capturing the cost the MTD imposes).
Hong et al.~\cite{hong2018} extend this partition with network-state-dynamic
measures (marked $\dagger$ in Table~\ref{tab:mtd-metrics}) that quantify how
the security posture shifts as MTD operations execute across consecutive
network states.

\begin{table*}[!t]
  \centering
  \caption{Representative MTD metrics by perspective and purpose.}
  \label{tab:mtd-metrics}
  \setlength{\tabcolsep}{10pt}
  \renewcommand{\arraystretch}{1.15}
  \begin{tabular}{@{}lll@{}}
    \toprule
    \textbf{Perspective} & \textbf{Effectiveness} & \textbf{Efficiency} \\
    \midrule
    Attacker-side &
      \begin{tabular}[t]{@{}l@{}}
        Attack success probability (ASP)~\cite{cho2020}\\
        Mean time to compromise (MTTC)~\cite{cho2020}\\
        Attack surface~\cite{cho2020}\\
        Attack path variation (APV)$^\dagger$~\cite{hong2018}\\
        Attack path number (APN)$^\dagger$~\cite{hong2018}\\
        Attack path exposure (APE)$^\dagger$~\cite{hong2018}
      \end{tabular} &
      \begin{tabular}[t]{@{}l@{}}
        Attack cost~\cite{cho2020}\\
        Payoff penalty~\cite{cho2020}
      \end{tabular} \\
    \addlinespace
    Defender-side &
      \begin{tabular}[t]{@{}l@{}}
        Defence success probability (DSP)~\cite{cho2020}\\
        Mean time to failure (MTTF)~\cite{cho2020}\\
        System security (CIA)~\cite{cho2020}\\
        Degree of vulnerability~\cite{cho2020}
      \end{tabular} &
      \begin{tabular}[t]{@{}l@{}}
        Defence cost~\cite{cho2020}\\
        Quality of service (QoS)~\cite{cho2020}\\
        System overhead~\cite{cho2020}\\
        Node variant cost (NVC)$^\dagger$~\cite{hong2018}\\
        Edge variation cost (EVC)$^\dagger$~\cite{hong2018}
      \end{tabular} \\
    \bottomrule
    \addlinespace[2pt]
    \multicolumn{3}{@{}l@{}}{\footnotesize\textit{Partition adapted from~\cite{cho2020}. $^\dagger$~network-state-dynamic measures, computed across consecutive network states.}} \\
  \end{tabular}
\end{table*}

With model and metrics specified, the remaining commitment is validation
method. Cho et al.~\cite{cho2020} classify these into four categories ---
analytical, simulation, emulation, and real testbed --- with the analytical
category splitting into probabilistic and graphical-security-model approaches.

The lineage this dissertation extends combines a graphical-security-model
structural representation with discrete-event simulation as the execution
engine. Brown et al.~\cite{brown2023} introduce MTDSim, a discrete-event
simulator built on a three-layer Hierarchical Attack Representation Model (HARM);
its principal
contribution is enabling combined-MTD evaluation, comparing shuffle,
diversity, and hybrid combinations against attack scenarios drawn from the
Cyber Kill Chain and MITRE ATT\&CK\@. MTDSim is the codebase this
dissertation extends; subsequent UWA work continues the line. The threat model itself --- named at the outset but examined nowhere above --- is the subject of Section~\ref{sec:convergence}, which argues it to be the dimension MTD evaluation most consistently leaves under-developed.



\subsection{MTD Orchestration and Adaptive Selection}
\label{subsec:mtd_orchestration}

Adaptive selection --- the choice of which MTD operation to deploy
at which moment --- is the frontier of defender sophistication, and
where the orchestration literature has concentrated. Across the
proactive, reactive, and hybrid scheduling regimes, the most capable
approaches sit between the extremes: pure proactive movement wastes
resources observing nothing, while pure reactive movement depends on
detection that is rarely complete or timely~\cite{cho2020}. How that
choice is made spans a spectrum from specified to learned --- explicit,
designer-authored rules at one end, a policy learned from interaction at
the other. The two works below mark its poles.

At the specified pole, the orchestration \emph{is} a set of explicit
conflict-resolution rules. Masud et al.~\cite{masud2025} address
the ``when to move'' question within a three-layer Temporal HARM (THARM) for IoT networks, where
MTD-mechanism coexistence is itself the orchestration problem: IP
shuffling and redundancy both modify network-layer attack paths and
cannot run concurrently, while OS diversity operates at the service
layer and can. Their trigger remains periodic and proactive, but
registration of an MTD operation enters a priority queue, and operations
blocked by an active conflict go to a suspension list until the
contested resource frees~\cite{masud2025}. The contribution lies in the conflict-resolution rules rather than the
trigger logic: these rules determine which mechanism deploys when
several are eligible. Evaluation across
attack risk, attack cost, and Return on Attack (RoA; a risk-to-cost ratio) shows combined
shuffle+diversity outperforming any single mechanism, at the cost of the
coexistence logic the orchestration layer must explicitly manage.

At the learned pole, specified rules give way to a policy learned from
interaction. Tay's UWA-lineage work~\cite{tay2024} sits
on the same MTDSim discrete-event platform examined in
Section~\ref{subsec:mtd_eval}~\cite{brown2023}, and trains a Double Deep
Q-Network (DDQN) agent that receives network security metrics at each
evaluation tick and selects one of five actions: IP shuffle, OS
diversity, service diversity, complete topology shuffle, or no
operation. The fifth action is structurally important --- restraint is
encoded as a policy choice rather than a global interval parameter,
which is the reinforcement learning (RL) paradigm's own answer to the ``when not to move'' half
of the lever.

Whether specified or learned, orchestration concentrates the defender's
growing sophistication on one side of the engagement. Tay's agent is one
example among a small but growing body of RL-orchestrated
multi-mechanism MTD; the broader paradigm inherits its training
adversary from whatever simulator it is built on. The defender's
decision-making has advanced; the adversary against which it is
evaluated is inherited, not modelled. That asymmetry is what the rest of
this review takes up.